% ============================================================
% Figure 5 melanoma clinical-framing block — Tumor Biology authorship
% Owner: Tumor Biology Agent (agent-mozurh1r)
% Co-owned with Immunology (lead, agent-mozus3vv), Data Science (figure, agent-mozutvsy), ML (pipeline, agent-mozurrmd)
% Target: sections/07_results_discovery.tex (Fig 5 main panel) + Fig 5 caption in main.tex
% Companion: Immunology's workspace/handoffs/fig5_tcf7_panels.md (agent/immunology@07fb8d4)
% Anchor: Immunology's workspace/fig5_tpex_substate_proposal.md
% ============================================================

% -------------------------------------------------------------
% INSERT AS A NEW SUBSECTION after the Fig 5 candidate-cluster discovery subsection
% (or as a standalone Fig 5 caption-block). Four-part structure: cohort table,
% ICB stratification, dose-response, negative control.
% -------------------------------------------------------------

\subsection{Melanoma clinical framing of the TCF7$^{\text{low}}$/SLAMF6$^{+}$ sub-state}
\label{sec:fig5-melanoma-clinical}

\textbf{Tumor Biology co-authorship; immunological sub-state calls from Immunology (see companion marker panels, SI Note S5).}

\paragraph{Cohort selection.}
We evaluate the TCF7$^{\text{low}}$/SLAMF6$^{+}$ sub-state on three melanoma cohorts with complementary properties (Table~\ref{tab:fig5-melanoma-cohorts}): the melanoma subset of the Kang et al.\ 2024 pan-cancer immune atlas \citep{Kang2024PanCancerImmune} as the scalability discovery substrate; Sade-Feldman et al.\ 2018 \citep{Sade-Feldman2018Cell} as the canonical longitudinal cohort (11 patients with matched pre-treatment and on-treatment anti-PD-1$\pm$anti-CTLA-4 samples); and Li et al.\ 2019 \citep{Li2019Cell} as the largest available cohort (32 patients spanning treatment-naive and anti-PD-1-treated, responder and non-responder). Wu et al.\ 2020 \citep{Wu2020NatCommun} (16 anti-PD-1-treated patients) is held for replication in Extended Data; its ICB-treated arm overlaps with Li et al.\ 2019.

\begin{table}[h]
\small
\caption{Melanoma cohorts used for Fig.\ 5 clinical framing of the TCF7$^{\text{low}}$/SLAMF6$^{+}$ sub-state.}
\label{tab:fig5-melanoma-cohorts}
\begin{tabular}{p{0.20\columnwidth}rp{0.18\columnwidth}p{0.32\columnwidth}}
\toprule
Cohort & $n$ patients & Treatment arm & Role \\
\midrule
Kang 2024 MEL & 87 & mixed & Scalability discovery substrate \\
Sade-Feldman 2018 & 11 & Anti-PD-1 $\pm$ CTLA-4, paired pre/on & Longitudinal primary \\
Li 2019 & 32 & Naive + anti-PD-1, R vs NR & Largest arm; outcome-association power \\
Wu 2020 & 16 & Anti-PD-1 & Held for Extended Data replication \\
\bottomrule
\end{tabular}
\end{table}

\paragraph{ICB-response stratification.}
The TCF7$^{\text{low}}$/SLAMF6$^{+}$ sub-state is identified using the marker panel in SI~Note~S5 \S2 (positive-core: SLAMF6, PDCD1$^{\text{mid}}$, CD8A, CXCR5$^{\text{variable}}$ [dim-to-negative accepted]; exclusion: TCF7$^{\text{hi}}$, LEF1$^{\text{hi}}$, CXCR5$^{\text{hi}}$, TOX$^{\text{hi}}$, ENTPD1, HAVCR2, GZMB, EOMES$^{\text{hi}}$). We define the per-patient sub-state frequency $f^{\text{sub}}_i$ as the fraction of CD8$^{+}$ TILs meeting the oracle criterion (Immunology SI Note S5 \S2 discriminating-test rule). Within each cohort, we test the hypothesis $f^{\text{sub}}_i$ differs between responders and non-responders under anti-PD-1 using a two-sided Mann-Whitney $U$-test with Benjamini-Hochberg correction across cohorts. We pre-register an effect-size threshold of $\log_2 \text{FC} \ge 0.5$ with $q < 0.05$ as the positive-discovery criterion, chosen to match published effect sizes for TPEX-vs-non-TPEX stratification on the same cohorts \citep{Sade-Feldman2018Cell,Miller2019NatImmunol}. Results will be reported as Panel~\ref{fig:f5a} forest plot across the three cohorts.

\paragraph{Dose-response within TPEX.}
We next test whether the sub-state frequency relative to canonical TCF7$^{\text{hi}}$/CXCR5$^{+}$ TPEX frequency is \emph{itself} a predictor of response, independent of total TPEX abundance. This is the dose-response claim: within patients who have comparable TPEX abundance, does the TCF7$^{\text{low}}$/SLAMF6$^{+}$ fraction shift the prediction? We compute a logistic-regression model of response on $\log(f^{\text{sub}}_i) - \log(f^{\text{TPEX}}_i)$, controlling for total TPEX abundance and cohort, and report the per-standard-deviation odds ratio in Panel~\ref{fig:f5c}. Controlling for total TPEX abundance addresses the reviewer concern that the sub-state could be a nuisance artefact of overall TPEX expansion under ICB rather than a distinct biomarker. A significant positive coefficient is the claim that REAL-detected sub-stratification within a canonical rare state carries biomarker information that the canonical state alone does not.

\paragraph{Negative control.}
The sub-state signal should be absent in cohorts not treated with checkpoint blockade. We test the same models on the Kang 2024 MEL subset restricted to non-ICB-treated samples (surgical resection, chemotherapy, targeted therapy arms) and report the null result in Panel~\ref{fig:f5d}. A strongly non-null coefficient here would indicate confounding by tumour-intrinsic biology, or by immune-status markers shared between the sub-state and disease severity; in that case the Fig.~5 ICB claim would be withdrawn and the sub-state would be demoted from biomarker-candidate to descriptive discovery (still significant for the REAL-detection claim but not for the clinical-framing one).

\paragraph{Pre-specified robustness.}
All four panels are repeated with (i) an alternate sub-state definition that relaxes the CXCR5$^{\text{variable}}$ constraint (per SI S5 Note on CXCR5 dim-to-negative), (ii) the hierarchical REAL scouting at the Tex-compartment ($\sim 100{,}000$ pooled CD8$^{+}$ cells) per Math Theorem~3 rather than whole-atlas scouting, and (iii) leave-one-cohort-out. Main-text figure reports the least-favourable robust estimate across the three sensitivity sweeps.

\paragraph{What would falsify the claim.}
The Fig.~5 ICB-biomarker claim is falsified if: (a) effect size below the pre-registered $\log_2 \text{FC} = 0.5$ in any two of three cohorts, (b) negative-control panel (Panel~\ref{fig:f5d}) shows a significant positive coefficient, or (c) the sub-state signal dissolves after controlling for total TPEX abundance in the dose-response panel. All three criteria are reported in the main figure regardless of outcome, per the evidence-card convention (SI~Note~S6; see \texttt{workspace/evidence\_cards/M-DEMO-TPEX-SUBSTATE.json} on agent/immunology).
